\documentclass[11pt]{article}
\usepackage[T1]{fontenc}
\usepackage{lmodern}
\usepackage{microtype}
\usepackage[margin=0.92in]{geometry}
\usepackage{amsmath,amssymb,amsthm,mathtools}
\usepackage{graphicx}
\usepackage{xcolor}
\usepackage{enumitem}
\usepackage[hidelinks]{hyperref}

\newtheorem{theorem}{Theorem}
\newtheorem{lemma}{Lemma}
\newcommand{\Hh}{\mathcal H}
\newcommand{\Rr}{\mathcal R}
\newcommand{\Kk}{\mathcal K}
\newcommand{\Dom}{\mathsf D}
\newcommand{\Ran}{\mathsf R}
\newcommand{\Ker}{\mathsf N}
\newcommand{\sgn}{\operatorname{sgn}}
\newcommand{\dist}{\operatorname{dist}}

\definecolor{statusblue}{RGB}{29,78,121}
\definecolor{softgray}{RGB}{245,247,249}

\title{Higher-Order Bounded Sign: Four Positive Regimes and an Abstract No-Go}
\author{Partial-result packet for arXiv:0810.3071}
\date{9 August 2026}

\begin{document}
\maketitle

\noindent\colorbox{softgray}{\parbox{0.95\linewidth}{
\textbf{Status.} \textcolor{statusblue}{Expanded partial result, likely valid
and requiring human review.} The source question is answered affirmatively
for fixed symbol range, for semidefinite $D$, when the compressed coefficient
is self-adjoint, and whenever the active compression is a one-sided spectral
multiplier perturbation $I+P_+M_bP_-$.  This last theorem allows rotating
range, both spectral signs, and non-self-adjoint coefficient. An abstract
counterexample shows why the full direction-dependent case cannot follow from
accretivity alone.
}}

\section*{Source question and framework}
Auscher, Axelsson, and McIntosh consider
$\Hh=L^2(\mathbb R^n;\mathbb C^m)$, a homogeneous constant-coefficient
differential operator $D$ of order $k$, and a bounded multiplication operator
$B$ \cite{AAM}. Their hypotheses (H1)--(H5) say that $D$ is self-adjoint,
\[
 \|\nabla^k u\|_2\lesssim\|Du\|_2
 \quad
 (u\in\Dom(D)\cap\overline{\Ran(D)}),
\]
and
\[
 \operatorname{Re}(BDu,Du)\geq\delta\|Du\|_2^2.
\]
On p.~2, immediately after Corollary 1.3, they ask whether for $k>1$
these hypotheses and the multiplication structure of $B$ already imply that
$\sgn(BD)$ is bounded. Their Theorem 1.1 proves the required quadratic
estimate after adding off-diagonal hypothesis (H6), and Proposition 1.2
deduces boundedness of $\sgn(BD)$ once that estimate is known for $B$ and
$B^*$.

\section*{Idea of the proof}
The obstruction in order $k>1$ is that a cutoff commutator $[D,\eta]$ contains
all derivative orders below $k$. If $D$ has dense range, however, (H3) is a
full elliptic estimate. The resolvent equation controls $t^kDu$, ellipticity
controls $t^k\nabla^ku$, and Fourier interpolation controls every
$t^j\nabla^ju$. Thus conjugating the resolvent by an exponential cutoff
changes it by a uniformly small operator from the scaled graph domain to
$L^2$. A Neumann series gives exponential off-diagonal decay, precisely (H6).

When the symbol range is a fixed subspace $V$, the kernel is the local fixed
space $V^\perp$. Compressing $B$ to $V$ preserves multiplication and
accretivity. The dense-range argument applies on $L^2(V)$, and elementary
similarities transfer the sign calculus back to $BD$.

For arbitrary symbol range, the same algebraic compression still works but is
nonlocal. It immediately settles two further regimes: a positive compression
is similar to a self-adjoint operator when it is itself self-adjoint, while a
semidefinite $D$ leaves only one spectral half-plane. Finally, an explicit
bilateral discrete-Hilbert-transform example makes the sign unbounded for an
abstract accretive compression with exponentially spaced spectrum. Thus a full
proof must exploit Euclidean polynomial growth and multiplier locality.

That polynomial growth can in fact be exploited for every one-sided spectral
compression $I+P_+M_bP_-$.  The spectral graph satisfies a linear Sylvester
equation.  Its Fourier kernel is split according to the relative sizes of the
input and output frequencies: comparable scales are Coifman--Meyer, the
high-output region is a standard high--low paraproduct, and the low-output
region gains a summable power of the scale ratio.  The resulting
$L^\infty\times L^2\to L^2$ bound gives an exact similarity to the
positive/negative diagonal operator.  The remaining general case is harder
precisely because two-sided coupling replaces this linear equation by a
nonlinear Riccati equation.

\begin{center}
\includegraphics[width=0.92\linewidth]{figures/open_problem_crop.png}

{\small The source proposition, first-order result, and higher-order question
(locally rendered arXiv:0810.3071v2, p.~2).}
\end{center}

\section*{Dense-range off-diagonal lemma}
\begin{lemma}\label{lem:dense}
Assume (H1)--(H5), let $B$ be a multiplication operator, and suppose
$\overline{\Ran(D)}=\Hh$. Then, for both choices of sign,
\[
 R_t^\pm=(I\pm i t^kBD)^{-1}
\]
has exponential $L^2$ off-diagonal decay at scale $t$. Consequently (H6)
holds for $B,D$, the same is true for $B^*,D$, and $\sgn(BD)$ is bounded.
\end{lemma}

\begin{proof}
Density extends (H5) to every $v\in\Hh$. Hence $B$ is strictly accretive on
$\Hh$ and has a bounded inverse. The source's bisectorial resolvent estimate
under (H2), (H4), and (H5) gives
\[
 \|u\|_2\lesssim\|(I\pm i t^kBD)u\|_2.
\]
Writing $L_t^\pm=I\pm i t^kBD$ and using $B^{-1}$ gives
\[
 t^k\|Du\|_2\lesssim\|L_t^\pm u\|_2.
\]
Since the range is dense, (H3) applies to every $u\in\Dom(D)$. Homogeneity
and the corresponding upper estimate show that $D$ is elliptic,
$\Dom(D)=H^k$, and
\[
 t^k\|\nabla^ku\|_2\lesssim\|L_t^\pm u\|_2.
\]
The elementary Fourier inequality
$(t|\xi|)^j\lesssim1+(t|\xi|)^k$, $0\leq j\leq k$, now yields the uniform
scaled graph estimate
\begin{equation}\label{eq:graph}
 \sum_{j=0}^k t^j\|\nabla^ju\|_2
 \lesssim \|L_t^\pm u\|_2.
\end{equation}

Let closed sets $E,F$ have distance $d>0$ and first suppose $t\leq d$.
Choose a smooth $\phi$ equal to $1$ on $E$ and $0$ on $F$, with
$\|\nabla^j\phi\|_\infty\lesssim d^{-j}$ for $1\leq j\leq k$. Put
$g=\alpha d\phi/t$, where $0<\alpha\leq1$ will be fixed. Leibniz' rule and
$t\leq d$ give
\[
 e^gDe^{-g}=D+K_g,
 \qquad
 t^k\|K_gu\|_2
 \lesssim \alpha\sum_{j=0}^{k-1}t^j\|\nabla^ju\|_2.
\]
Indeed, a term in which $j$ derivatives hit $e^{-g}$ has coefficient at most
$C\alpha t^{-j}$; products of several derivatives of $g$ only add powers of
$\alpha\leq1$. Combining this with \eqref{eq:graph} shows
\[
 \|i t^kB K_g(L_t^\pm)^{-1}\|_{2\to2}\leq C\alpha.
\]
Choose $\alpha$ so that $C\alpha<1/2$. Then
\[
 e^gL_t^\pm e^{-g}=L_t^\pm\pm i t^kB K_g
\]
is invertible with a uniform $L^2$ inverse by a Neumann series. If
$f$ is supported in $F$, then $e^gf=f$, whereas $e^g=e^{\alpha d/t}$ on
$E$. Therefore
\[
 \|1_E R_t^\pm f\|_2\lesssim e^{-\alpha d/t}\|f\|_2.
\]
For $t>d$ the uniform resolvent bound gives the same estimate after changing
the constant. Finally,
\[
 t^kBD(I+t^{2k}BDBD)^{-1}
   =\frac{1}{2i}(R_t^- -R_t^+),
\]
so (H6) follows, with exponential rather than merely polynomial decay. The
same proof applies to $B^*$ because it has the same accretivity real part.
Theorem 1.1 and Proposition 1.2 of \cite{AAM} now imply boundedness of
$\sgn(BD)$.
\end{proof}

\section*{Active-space reduction}
Let $P$ be the orthogonal projection onto
$\Rr=\overline{\Ran(D)}$, let $\Kk=\Ker(D)$, and write $D_R$ for the
restriction of $D$ to $\Rr$.

\begin{lemma}\label{lem:reduction}
Put $C=P B|_{\Rr}$ and $U=B|_{\Rr}$. Then $C$ is boundedly invertible on
$\Rr$, $U$ is an isomorphism from $\Rr$ onto $B\Rr$, and
\[
 \Hh=B\Rr\dotplus\Kk,\qquad BD\,U=U D_R C.
\]
Consequently boundedness of $\sgn(BD)$ is equivalent to boundedness of
$\sgn(D_R C)$.
\end{lemma}

\begin{proof}
By (H5) and density of $\Ran(D)$ in $\Rr$,
$\operatorname{Re}(Cv,v)\geq\delta\|v\|^2$ on $\Rr$, so $C$ is invertible
and $U$ is bounded below. For $h\in\Hh$, set $v=C^{-1}Ph$. Then
$h-Bv\in\Kk$, and the sum is direct because $PBv=Cv$. Finally,
$DBv=DPBv=D_RCv$, which gives the intertwining identity. Functional calculus
is invariant under the displayed bounded similarity and is zero on $\Kk$.
\end{proof}

\section*{Three affirmative regimes}
Let $\widehat D(\xi)$ denote the homogeneous matrix symbol of $D$.

\begin{theorem}\label{thm:fixed}
Assume (H1)--(H5), let $B$ be a bounded multiplication operator, and suppose
there is a fixed subspace $V\subseteq\mathbb C^m$ such that
\[
 \Ran\widehat D(\xi)=V \qquad(\xi\ne0).
\]
Then $\sgn(BD)$ is bounded on $L^2(\mathbb R^n;\mathbb C^m)$ and invertible
on $\overline{\Ran(BD)}$.
\end{theorem}

\begin{proof}
Here $P$ is the constant pointwise projection onto $V$, so
$C=P B|_V$ is again a multiplication operator. The operator $D_R$ has dense
range on $L^2(V)$, and Lemma~\ref{lem:dense} gives bounded $\sgn(CD_R)$.
Since $D_RC=C^{-1}(CD_R)C$, Lemma~\ref{lem:reduction} transfers the bounded
sign to $BD$.
\end{proof}

\begin{theorem}\label{thm:two-more}
Under (H1)--(H5), $\sgn(BD)$ is bounded in either of the following cases:
\begin{enumerate}[label=(\roman*)]
\item $D$ is nonnegative or nonpositive;
\item $C=P B|_{\Rr}$ is self-adjoint on $\Rr$; in particular, this holds
when $B$ is self-adjoint on $\Hh$.
\end{enumerate}
No fixed-range hypothesis is required.
\end{theorem}

\begin{proof}
If $D_R\geq0$, then $CD_R$ is sectorial in the right half-plane, by strict
accretivity of $C$, and its sign on $\Rr$ is $I$. If $D_R\leq0$, its sign is
$-I$. Since $D_RC=C^{-1}(CD_R)C$, conclusion (i) follows from
Lemma~\ref{lem:reduction}.

For (ii), (H5) gives $C\geq\delta I$ on $\Rr$. The operator
\[
 A=C^{1/2}D_RC^{1/2},
 \qquad \Dom(A)=C^{-1/2}\Dom(D_R),
\]
is self-adjoint, and
$D_RC=C^{-1/2}AC^{1/2}$. Thus $\sgn(D_RC)$ is bounded, and the reduction
again finishes the proof.
\end{proof}

\section*{A general one-sided spectral-triangular theorem}
Let $P_+$ and $P_-$ be the positive and negative spectral projections of
$D_R$, and write $\Rr_\pm=P_\pm\Rr$.  Hypothesis (H3) gives a uniform gap
between zero and the nonzero eigenvalues of the homogeneous symbol.  Thus
$P_\pm$ are smooth degree-zero matrix multipliers away from the origin, and
\[
 D_R=\begin{bmatrix}A_+&0\\0&-A_-\end{bmatrix}
 \quad\hbox{on }\Rr_+\oplus\Rr_-,
\]
where $A_\pm$ are positive bundle multipliers whose eigenvalues are comparable
to $|\xi|^k$.

\begin{lemma}[Bundle Sylvester estimate]\label{lem:sylvester}
For $b\in L^\infty(\mathbb R^n;\mathcal L(\mathbb C^m))$, set
\[
 N=P_+M_bP_-\colon\Rr_-\longrightarrow\Rr_+.
\]
There is a bounded $X\colon\Rr_-\to\Rr_+$ such that
\begin{equation}\label{eq:sylvester}
 A_+X+XA_-=2A_+N,
 \qquad \|X\|\lesssim_D\|b\|_\infty.
\end{equation}
\end{lemma}

\begin{proof}
Let $a_+(\zeta)$ and $a_-(\eta)$ denote the positive bundle symbols.  For a
matrix $H$, let $\mathfrak m(\zeta,\eta)[H]$ be the unique bundle map $x$
solving
\begin{equation}\label{eq:fibresyl}
 a_+(\zeta)x+xa_-(\eta)
 =2a_+(\zeta)p_+(\zeta)H p_-(\eta).
\end{equation}
The positive Sylvester map on the left has inverse norm at most
$C_D/(|\zeta|^k+|\eta|^k)$.  Define $X$ by the bilinear Fourier multiplier
whose kernel is \eqref{eq:fibresyl}, with
$H=\widehat b(\zeta-\eta)$.  We prove its endpoint bound by a smooth partition
into three frequency-ratio regions.  The separation constant is chosen small
enough that the two Neumann series below converge in every finite collection
of multiplier seminorms; polynomial factors produced by differentiating the
$j$th term are absorbed by its geometric ratio gain.

When $|\zeta|\simeq|\eta|$, the normalized bundle symbol is smooth and
homogeneous of degree zero.  After factoring off the bounded input and output
multipliers $P_-$ and $P_+$, it is an ordinary bilinear Coifman--Meyer symbol.
The standard endpoint estimate $L^\infty\times L^2\to L^2$ therefore applies
\cite[Introduction]{ARL}.

Suppose next that $|\eta|\ll|\zeta|$ and put
$K=p_+(\zeta)H p_-(\eta)$.  Multiplying \eqref{eq:fibresyl} on the left by
$a_+(\zeta)^{-1}$ gives the norm-convergent expansion
\begin{equation}\label{eq:highlow}
 x=2\sum_{j\geq0}(-1)^j
 a_+(\zeta)^{-j}K a_-(\eta)^j.
\end{equation}
The $j=0$ term is the standard high--low paraproduct; its
$BMO\times L^2\to L^2$ bound applies because $L^\infty\subset BMO$.  On each
dyadic ratio layer, the $j$th remaining term gains
$(|\eta|/|\zeta|)^{kj}$.  The normalized angular multipliers are uniformly
Calder\'on--Zygmund, so the paraproduct estimates and the geometric series in
$j$ and in the ratio layers give the required bound.

Finally, when $|\zeta|\ll|\eta|$, multiplication on the right by
$a_-(\eta)^{-1}$ gives
\begin{equation}\label{eq:lowoutput}
 x=2\sum_{j\geq0}(-1)^j
 a_+(\zeta)^{j+1}K a_-(\eta)^{-j-1}.
\end{equation}
Every term now gains
$(|\zeta|/|\eta|)^{k(j+1)}$.  On the layer
$|\zeta|\simeq2^{-\ell}|\eta|$, the output projection of the product is
bounded by
\[
 C_D2^{-k\ell}\|b\|_\infty\|v_\eta\|_2.
\]
Littlewood--Paley orthogonality sums the input scales and the geometric
$\ell$-sum closes.  Equations \eqref{eq:highlow} and
\eqref{eq:lowoutput}, together with the comparable region, prove
$\|X\|\lesssim_D\|b\|_\infty$.  Equation \eqref{eq:sylvester} follows from
\eqref{eq:fibresyl} in distributions.
\end{proof}

\begin{theorem}\label{thm:triangular}
Assume (H1)--(H5).  If, relative to
$\Rr=\Rr_+\oplus\Rr_-$, the active compression has the form
\[
 C=P B|_{\Rr}=
 \begin{bmatrix}I&P_+M_bP_-\\0&I\end{bmatrix}
\]
for some bounded matrix function $b$, then $\sgn(BD)$ is bounded.  The same
holds with the lower-triangular orientation reversed.
\end{theorem}

\begin{proof}
Let $X$ be given by Lemma~\ref{lem:sylvester} and put
\[
 R=\begin{bmatrix}I&-X/2\\0&I\end{bmatrix},
 \qquad J=\begin{bmatrix}A_+&0\\0&-A_-\end{bmatrix}.
\]
The Sylvester identity says
$A_+(N-X/2)=XA_-/2$.  Hence $N-X/2$ maps
$\Dom(A_-)$ into $\Dom(A_+)$, and a direct calculation, including the reverse
domain inclusion, gives the closed-operator identity
\[
 D_RC=RJR^{-1}.
\]
Functional calculus under this bounded similarity yields
\[
 \sgn(D_RC)
 =R\begin{bmatrix}I&0\\0&-I\end{bmatrix}R^{-1}
 =\begin{bmatrix}I&X\\0&-I\end{bmatrix}.
\]
Lemma~\ref{lem:reduction} transfers the bounded sign to $BD$.  Taking
adjoints, or reversing the two bundles, proves the lower-triangular case.
\end{proof}

For example, Theorem~\ref{thm:triangular} contains the following genuinely
rotating-range family in every even order.  On
$L^2(\mathbb R^n;\mathbb C\oplus\mathbb C^n)$ let
\[
 \widehat D_r(\xi)=
 \begin{bmatrix}|\xi|^{2r}&0\\
 0&-|\xi|^{2r-2}\xi\xi^*\end{bmatrix},
 \qquad
 B_\beta(x)=\begin{bmatrix}1&\beta(x)\ell\\0&I_n\end{bmatrix}.
\]
Here the negative range is the radial line and the kernel is the tangential
plane.  If $\|\beta\|_\infty\|\ell\|<2$, then $B_\beta$ is strictly
accretive and the theorem applies.  In this case the Sylvester symbol reduces
to the explicit Coifman--Meyer function
\[
 \frac{2|\alpha+\eta|^{2r}}
 {|\alpha+\eta|^{2r}+|\eta|^{2r}}.
\]

\section*{Why accretivity alone is insufficient}
The next example is not a Euclidean differential counterexample. It shows
that the locality and polynomial spectral growth in the source problem are
indispensable.

\begin{theorem}[Abstract no-go]\label{thm:nogo}
There are an injective self-adjoint operator $D_0$ with dense range and a
bounded uniformly strictly accretive operator $C_0$ on a Hilbert space for
which $\sgn(C_0D_0)$ is unbounded.
\end{theorem}

\begin{proof}
Let $H_0=\ell^2(\mathbb Z)$, fix $r>1$, and set
$Ae_j=r^je_j$. Let $T$ be bilateral convolution with
$h_d=1/d$ for $d\neq0$ and $h_0=0$. Its Fourier symbol is the bounded
sawtooth series $\sum_{d\neq0}e^{id\theta}/d$, so $T$ is bounded. On
$H_0\oplus H_0$ put
\[
 D_0=\begin{bmatrix}A&0\\0&-A\end{bmatrix},\qquad
 C_0=\begin{bmatrix}I&\varepsilon T\\0&I\end{bmatrix},
 \qquad 0<\varepsilon\|T\|<1.
\]
The Hermitian part of $C_0$ is bounded below by
$1-\varepsilon\|T\|/2>0$.

If $\sgn(C_0D_0)$ were bounded, triangular functional calculus would give
\[
 \sgn(C_0D_0)=\begin{bmatrix}I&X\\0&-I\end{bmatrix},
 \qquad AX+XA=-2\varepsilon TA.
\]
Matrix entries force, for $i\neq j$,
\[
 X_{ij}=\frac{-2\varepsilon}{(i-j)(1+r^{i-j})}.
\]
Thus $X$ would be convolution by
$x_d=-2\varepsilon/[d(1+r^d)]$. For $d\to-\infty$ this has the one-sided
harmonic tail $-2\varepsilon/d$, while for $d\to+\infty$ it decays
exponentially. Its Fourier series is therefore a nonzero multiple of
$\log(1-e^{-i\theta})$ plus a bounded function, and is not essentially
bounded. Such convolution is unbounded on $\ell^2(\mathbb Z)$, a
contradiction.
\end{proof}

\section*{Consequences and remaining obstruction}
The hypothesis of Theorem~\ref{thm:fixed} holds whenever $D$ is elliptic
($V=\mathbb C^m$), whenever $D$ is scalar and nonzero under (H3), and more
generally whenever the kernel of $\widehat D(\xi)$ is independent of
$\xi\ne0$.

Theorem~\ref{thm:two-more} also settles every semidefinite system and every
self-adjoint coefficient, even when the symbol range rotates.
Theorem~\ref{thm:triangular} goes further for every homogeneous symbol $D$:
it settles all one-sided spectral compressions $I+P_+M_bP_-$ and their
lower-triangular analogues, including non-self-adjoint examples with both
signs and rotating range.
The unresolved core therefore requires genuinely two-sided coupling between
the positive and negative spectral bundles; having both signs, rotating
range, and non-self-adjoint compression is no longer by itself an obstruction.

This is not merely a missing estimate in the proof above. If
$\Ran\widehat D(\xi)$ varies with direction, the projection onto
$\overline{\Ran(D)}$ is a nonlocal homogeneous Fourier multiplier. Compressing
$B$ to that range therefore destroys the local multiplication structure, and
exponential conjugation of the projection is not uniformly bounded.
Theorem~\ref{thm:nogo} shows that accretivity alone cannot replace this lost
locality. A full solution must exploit polynomial homogeneous spectrum and
the fact that $B$ is multiplication, or construct a Euclidean example that
defeats those structures.  For a general block compression, eliminating one
off-diagonal block introduces inverses of nonlocal Toeplitz-type compressions;
equivalently the linear Sylvester equation in
Theorem~\ref{thm:triangular} becomes a nonlinear Riccati equation without a
contraction parameter.

\section*{Verification and novelty check}
The graph estimate tracks every power of $t$ and derivative order in the
higher-order commutator; its constant is independent of $t,d,E,F$. The
argument was checked separately for $B^*$, as required by source Proposition
1.2. The active-space reduction was checked without assuming that the range
projection is local. In the no-go example, the forced off-diagonal block is a
convolution operator and its unbounded logarithmic Fourier symbol gives an
exact contradiction, not numerical evidence.

For Lemma~\ref{lem:sylvester}, the audit separates all three frequency-ratio
regions.  Comparable scales use the cited endpoint Coifman--Meyer estimate;
the high-output region is the standard $BMO\times L^2$ paraproduct plus a
geometric Sylvester series; and the low-output region has the explicit
$2^{-k\ell}$ gain needed to sum its ratio layers.  The fibrewise Sylvester
identity then gives the operator identity in distributions.  Finally,
$A_+(N-X/2)=XA_-/2$ supplies the exact domain mapping in the triangular
similarity.  The displayed even-order family independently checks the formula
against its scalar explicit symbol.

No numerical assertion enters the proof. Finite Fourier experiments were
discarded because they produced the same apparent growth in the known
first-order case. The run's four cheap indexes had no hit for the source id or
theorem. Bounded local citation scans and web searches on 9 August 2026 for the
exact question, higher-order $BD$ calculus, dense range, fixed kernel, and
fixed symbol range found the source and adjacent first-order/structured Kato
work.  Additional exact-phrase searches for triangular coefficients,
spectral compressions $P_+M_bP_-$, rotating-range signs, and the
Sylvester/Coifman--Meyer mechanism found the source and general paraproduct
literature, but no explicit statement of
Theorem~\ref{thm:triangular}. Novelty confidence is moderate because the
literature check was bounded and the area is specialized.

\textbf{Human-review recommendation.} Check the Leibniz bound for $K_g$ in
Lemma~\ref{lem:dense}, the unbounded one-sided harmonic Fourier series in
Theorem~\ref{thm:nogo}, and the domain statement for
$C^{1/2}D_RC^{1/2}$ in Theorem~\ref{thm:two-more}.  For the new theorem,
verify the three-region dyadic estimate in Lemma~\ref{lem:sylvester}, with
particular attention to the $BMO$ paraproduct at $j=0$, and the domain passage
in the triangular similarity.

\begin{thebibliography}{9}
\bibitem{AAM}
P. Auscher, A. Axelsson, and A. McIntosh,
\emph{On a quadratic estimate related to the Kato conjecture and boundary
value problems}, Harmonic Analysis and Partial Differential Equations,
Contemp. Math. \textbf{505} (2010), 105--129; arXiv:0810.3071v2.

\bibitem{ARL}
S. Arias and S. Rodr\'iguez-L\'opez,
\emph{Some endpoint estimates for bilinear Coifman--Meyer multipliers},
arXiv:2005.11771v1 (2020).
\end{thebibliography}

\end{document}
